% !TeX root = ../教材题库_全部条目_转写工作稿.tex
% Source: XJ-B1-C06-EX-P251-Y1103
% 第6章 / 6.4 用样本估计总体 / 教材 p.251、252 / PDF 第258、259页
% OCR未校对：公式与文字均可能有误；未完成数学TeX转写；禁止直接发给学生
% 原题及相邻原解截图（核对用）：../题目原图/XJ-B1-C06-EX-P251-Y1103-01.webp, ../题目原图/XJ-B1-C06-EX-P251-Y1103-02.webp
% 本文件只在转写完成后作为题干使用；OCR 条目尚不是可靠题干。
\begin{bankitem}{6.4 例题 例题10}{第6章 / 6.4 用样本估计总体 / 教材 p.251、252 / PDF 第258、259页}{XJ-B1-C06-EX-P251-Y1103}{OCR未校对：公式与文字均可能有误；未完成数学TeX转写；禁止直接发给学生}
\begin{ocrdraft}
[教材p.251 / PDF第258页]
\par
0：，0]计算数据
\par
1，
\par
5，9，
\par
12，
\par
13，
\par
18，
\par
21，
\par
23，
\par
28，
\par
36
\par
的百分位数：P25，P
\par
，P75
\par
解数据量7冖10．
\par
因为一\ensuremath{\times}25％一2．5不是整数，3是比2．5大的最小整数，
\par
所以P25一一9；
\par
因为c=n><50％一5是整数，
\par
5\ocrglyph{00AF}f\ocrglyph{00AF}X6
\par
一13十18
\par
所以P
\par
50
\par
=15．5；
\par
2
\par
2
\par
因为一\ensuremath{\times}75％一7．5不是整数，8是比7．5大的最小整数，
\par
所以P7一“一23．
\par
对于从小到大排列的个数，P2
\par
大约处于这77个数的了处，P50大约处于这77个
\par
[教材p.252 / PDF第259页]
\par
数的中间，P7大约处于这个数的一处．
\par
百分位数是用于衡量数据位置的度量，它提供了有关数据在最小值与最大值之
\par
间位置的信息．多个百分位数结合使用，可更全面地描述数据的分布特征．
\end{ocrdraft}
\end{bankitem}
